% ============================================================
% Appendix Z.6 — FRFP Case Study (Law / Professional Practice):
% Mata v. Avianca and the Hallucinated Citations
% ============================================================

\section{FRFP Case Study (Law / Professional Practice): Mata v.\ Avianca and the Hallucinated Citations}
\label{app:frfp-mata-avianca}

This appendix examines \emph{Mata v.\ Avianca, Inc.} and the resulting
sanctions for AI-generated fake case citations as a case study in FRFP terms.
We treat the episode not as an isolated lawyer error, but as an instance of
explicit/tacit confusion, boundary failure, and absent safe horizons in
professional use of AI tools.

\subsection{Background}

In \emph{Mata v.\ Avianca, Inc.}, a personal injury suit in the Southern
District of New York, plaintiff’s counsel used ChatGPT to generate portions
of a brief opposing a motion to dismiss.
The brief contained citations to multiple cases that did not exist, complete
with plausible but fabricated quotations and docket numbers.

When Avianca’s lawyers and the court were unable to locate several cited
cases, Judge P.\ Kevin Castel ordered plaintiff’s counsel to provide copies.
Counsel again turned to ChatGPT, which falsely asserted that the cases were
real and available in standard databases.
The lawyers then filed AI-generated ``opinions'' as if they were genuine.

After an evidentiary hearing, Judge Castel sanctioned the two attorneys and
their firm under Rule~11 and the court’s inherent authority, imposing a
\$5{,}000 fine and requiring notification of the misconduct to affected
parties.
His opinion emphasised that:

\begin{itemize}
  \item the attorneys filed non-existent cases without adequate verification;
  \item they persisted after warning signs (defence counsel and the court
        could not locate the cases);
  \item they misrepresented the status of their research to the court.
\end{itemize}

The case quickly became a touchstone for AI use in legal practice.
Subsequent commentary and ethics opinions now routinely cite \emph{Mata v.\
Avianca} as a landmark on lawyers’ duties when using generative AI.

\subsection{Explicit and Tacit Spaces in Legal Practice}

\subsubsection*{Explicit space \(E_{\mathrm{law}}\)}

In FRFP terms, the explicit domain includes:

\begin{itemize}
  \item legal texts: statutes, regulations, case law, treatises, and
        secondary sources;
  \item legal documents: complaints, motions, briefs, affidavits, and
        exhibits;
  \item citations, quotation blocks, and argument structures in filings;
  \item AI outputs: draft arguments, case lists, summaries, and suggested
        citations.
\end{itemize}

All of this is manipulable by AI systems.
Given prompts, a model can generate briefs, summaries, and case references
that appear syntactically and stylistically plausible.

\subsubsection*{Tacit space \(T_{\mathrm{law}}\)}

Tacit space contains:

\begin{itemize}
  \item lawyers’ understanding of doctrinal structure, relevance, and
        persuasive strategy;
  \item professional judgment about what is arguable, what is frivolous,
        and what is ethical;
  \item local knowledge of court norms, opposing counsel, and practical
        litigation constraints;
  \item ethical commitments and fear of sanctions or reputational harm.
\end{itemize}

Correctness in law is not just about textual fit:
it involves tacit assessments of authority, applicability, and integrity.

\subsection{Boundary Design: Filing as a Tacit Endorsement}

\subsubsection*{Formal boundary in litigation practice}

Under FRFP, each filing is a boundary artefact:

\begin{itemize}
  \item a lawyer’s signature (physical or electronic) on a brief or motion is
        a boundary event in \(T_{\mathrm{law}}\);
  \item by signing, counsel tacitly endorses that:
        \begin{itemize}
          \item factual content has evidentiary support;
          \item legal arguments are grounded in existing law or plausible
                non-frivolous extensions;
          \item citations correspond to real, relevant authorities.
        \end{itemize}
\end{itemize}

Rule~11 and analogous ethical rules explicitly codify this boundary:
they assign to lawyers the responsibility to verify the explicit artefacts
they submit, regardless of who or what drafted them.

\subsubsection*{Boundary confusion in Mata}

In \emph{Mata v.\ Avianca}, the boundary between AI output and lawyer
endorsement was blurred:

\begin{itemize}
  \item counsel treated ChatGPT’s outputs as if they were akin to traditional
        research tool results, assuming it would not fabricate cases;
  \item they did not independently verify the case citations in conventional
        databases before filing;
  \item even after opposing counsel and the court raised concerns, they
        returned to ChatGPT to ``confirm'' the existence of the cases rather
        than re-grounding in authoritative sources.
\end{itemize}

From an FRFP perspective:

\begin{itemize}
  \item the filing of the brief was treated as a mostly explicit operation
        (exporting and submitting AI text), rather than a tacit endorsement
        requiring intensive human evaluation;
  \item correctness authority was informally ceded to the tool, contrary to
        the formal boundary encoded in professional rules.
\end{itemize}

\subsection{Protocol-Level Failures in FRFP Terms}

\subsubsection*{Z.6.1 Explicit-Space Hallucination and Semantic Fragility}

\paragraph{Problem.}
Generative models are known to \emph{hallucinate}:
they produce fluent but fabricated content, including nonexistent cases and
fake citations.
In \emph{Mata}, ChatGPT:

\begin{itemize}
  \item generated case names, docket numbers, and quotes that mimicked real
        federal decisions;
  \item insisted, when queried, that the cases ``indeed exist'' and could be
        found in Lexis or Westlaw;
  \item produced PDFs with plausible formatting that further deceived counsel.
\end{itemize}

In FRFP language, the explicit semantics of the AI system for the task
``produce supporting cases'' is structurally unsafe:

\begin{itemize}
  \item the mapping from prompts to outputs does not preserve a key invariant
        (truthfulness or referential validity of citations);
  \item the system has no internal notion of ``being grounded in the corpus
        of real cases''; it only optimises for textual plausibility.
\end{itemize}

\paragraph{Consequence.}
Even without any negligence at the boundary, such a system \emph{cannot} be
treated as an oracle of case law.
It is, at best, a suggestion engine.
Treating its outputs as ready-to-file legal authority is an explicit-space
misinterpretation.

\subsubsection*{Z.6.2 Boundary Failure: From Tool to Surrogate Researcher}

\paragraph{Problem.}
The lawyers in \emph{Mata} effectively used ChatGPT as a surrogate legal
researcher:

\begin{itemize}
  \item they asked it for cases supporting their proposition;
  \item they accepted its answers without verifying them in standard research
        tools;
  \item they relied on its reassurances even after doubts were raised.
\end{itemize}

In FRFP terms:

\begin{itemize}
  \item the boundary between explicit AI generation and tacit legal
        endorsement collapsed: the AI’s text crossed almost directly into
        the filed brief;
  \item the human role reduced to light editing and formatting, not
        substantive verification;
  \item the signature, a boundary marker, became decoupled from actual
        tacit evaluation.
\end{itemize}

\paragraph{Consequence.}
This is a textbook instance of \emph{boundary capture}:
what should have been a human correctness function at filing time was, in
effect, automated by an explicit system whose semantics do not support that
role.

\subsubsection*{Z.6.3 Missing Safe Horizons and Re-Grounding}

\paragraph{Problem.}
Use of AI tools for legal drafting and research is inherently long-horizon:

\begin{itemize}
  \item lawyers may come to rely on such tools daily or hourly across cases;
  \item tacit skills in manual research and critical reading can degrade;
  \item oversight can become perfunctory as trust in the tool grows.
\end{itemize}

In \emph{Mata}, the episode occurred relatively early in the legal
profession’s encounter with generative AI, but the pattern is instructive:

\begin{itemize}
  \item counsel used AI extensively in one brief without clear constraints;
  \item there was no firm-level safe-horizon policy (e.g.\ limits on unvetted
        AI usage, mandatory verification steps);
  \item when the tool’s reliability was questioned, they did not reset to
        a ``manual'' mode, but deepened reliance on the same explicit system.
\end{itemize}

\paragraph{Consequence.}
FRFP would say that the run---the sequence of research and drafting steps
for that brief---had an unsafe horizon:
the depth of AI reliance exceeded the robustness of explicit semantics and
the available tacit oversight.
A well-designed FRFP implementation in law would enforce shorter horizons
and mandatory re-grounding in primary sources.

\subsection{Beyond Mata: A Pattern of Legal AI Failures}

\emph{Mata v.\ Avianca} was the first widely publicised sanction for AI
hallucinations in legal practice, but subsequent cases show it is not an
isolated event:

\begin{itemize}
  \item A Utah appeals court sanctioned an attorney in 2025 for filing a brief
        with fabricated citations generated by ChatGPT, including a nonexistent
        case, and ordered restitution and charitable donations.
  \item California’s Court of Appeal fined a lawyer \$10{,}000 after
        discovering that 21 of 23 quoted case passages in a brief were AI
        fabrications, noting the growing need for AI-use regulation.
  \item Multiple federal and state judges have now sanctioned or reprimanded
        lawyers and experts for unverified AI-generated citations, with
        surveys and bar commentary describing an ``epidemic'' of such
        errors.
\end{itemize}

These events illustrate a population-level FRFP concern:

\begin{itemize}
  \item as more practitioners incorporate AI into workflows,
        tacit professional norms about verification can degrade;
  \item explicit-only governance (tool warnings, EULA disclaimers) is
        insufficient to prevent structural misuse;
  \item courts are being forced, case by case, to reassert the boundary:
        correctness and responsibility remain with human lawyers.
\end{itemize}

\subsection{Counterfactual: An FRFP-Aligned Legal Research Workflow}

We sketch how FRFP might structure AI use in legal practice.

\subsubsection*{Z.6* Explicit/Tacit Role Separation}

\begin{itemize}
  \item AI tools are designated as \emph{drafting and ideation assistants}:
        they may propose arguments, case lists, and language.
  \item Human lawyers remain solely responsible for:
        \begin{itemize}
          \item selecting authorities to cite;
          \item verifying that citations correspond to real, relevant cases;
          \item assessing the doctrinal and strategic soundness of arguments.
        \end{itemize}
\end{itemize}

\subsubsection*{Z.6* Boundary Design at Filing}

\begin{itemize}
  \item Filing a document is treated as a formal boundary event:
        \begin{itemize}
          \item the lawyer must attest (internally, and where required by
                rule) that all citations have been verified using trusted
                sources (official reporters, recognised databases);
          \item firms implement checklists or tooling to mark each citation
                as ``verified'' before filing.
        \end{itemize}
  \item AI-generated content is visible in internal logs (e.g.\ authorship
        annotations in the draft) but not privileged in correctness:
        the final filing is a human artefact.
\end{itemize}

\subsubsection*{Z.6* Safe Horizons and Firm-Level Policy}

\begin{itemize}
  \item Firms and courts define safe-horizon policies for AI use:
        \begin{itemize}
          \item no filing may rely solely on AI-driven research; at least one
                human must independently run conventional research;
          \item heavy AI use in earlier drafts must be balanced by more
                intensive human verification at later stages;
          \item repeated use of AI tools triggers periodic audits of workflows
                and error rates.
        \end{itemize}
  \item Training and onboarding emphasise the limitations of generative
        models, explicitly correcting the misconception that they cannot
        ``fabricate'' content.
\end{itemize}

\subsubsection*{Z.6* Institutional Governance}

\begin{itemize}
  \item Bar associations and courts issue FRFP-style guidance:
        \begin{itemize}
          \item clarifying that AI assistance does not change the locus of
                professional responsibility;
          \item recommending or requiring disclosure of AI use in certain
                contexts;
          \item defining sanctions and remedial steps for AI-related
                misconduct.
        \end{itemize}
  \item Firms create internal AI oversight committees to:
        \begin{itemize}
          \item approve tools and acceptable use cases;
          \item monitor incidents and near-misses;
          \item update safe-horizon and verification procedures.
        \end{itemize}
\end{itemize}

\subsection{Lessons for FRFP and Professional Domains}

The \emph{Mata} case and its successors support several FRFP themes:

\begin{enumerate}
  \item \textbf{Generative AI is structurally hallucination-prone.}
        Its explicit semantics optimise for plausibility, not truth.
        Without structural constraints, it will produce fabricated but
        authoritative-looking artefacts.

  \item \textbf{Signatures are boundaries, not formalities.}
        In law and other professions, signing a document is a boundary event
        at which tacit endorsement is supposed to be applied.
        Automating the path up to that boundary does not diminish the
        responsibility that attaches at it.

  \item \textbf{Safe horizons are essential in professional workflows.}
        Long-horizon, high-frequency reliance on AI without explicit
        verification protocols invites degradation of tacit skills and norms.

  \item \textbf{Explicit-only governance is insufficient.}
        Tool warnings, footnotes, and vendor caveats cannot substitute for
        institutional operators (courts, bars, firms) that enforce FRFP-style
        boundaries and safe horizons.

  \item \textbf{Formal architectures clarify ``AI mishaps'' as structural.}
        \emph{Mata} is often narrated as a cautionary tale about a few careless
        lawyers.
        FRFP reframes it as a prototypical failure mode whenever explicit
        AI systems are allowed to leak across professional boundaries without
        preserving tacit correctness authority.
\end{enumerate}

As generative AI becomes embedded across professions, the lesson of
\emph{Mata v.\ Avianca} is that AI assistance can coexist with professional
responsibility only under architectures that explicitly maintain the
explicit/tacit split, design robust boundaries at endorsement points, and
enforce finite safe horizons with institutional oversight.
